
\subsubsection{6.1 Key Findings}

\textbf{Non-monotonic dose-response relationships are measurable.} The cubic fit ($R^2$ = 0.985) decisively outperforms linear models, confirming that curation parameter effects are not monotonic. This challenges the assumption in some pipelines that stricter filtering is uniformly better.

\textbf{The mechanism is noise dilution at low thresholds, diversity loss at high thresholds.} Convergence dynamics analysis shows 40\% higher AUC for unfiltered training, demonstrating that noisy data dilutes learning signals. Simultaneously, aggressive filtering achieves reasonable convergence but poor final performance, indicating loss of valuable information for generalization.

\textbf{Optimal thresholds show consistency across scales.} The transfer ratio of 0.85 between 125M and 1B models suggests small-scale sweeps can guide large-scale training with modest correction.

\subsubsection{6.2 Limitations}

\textbf{Reduced-scale evaluation.} All experiments used 5M--10M tokens rather than planned 10B, with mock/synthetic benchmark evaluation. Effect magnitudes require full-scale validation.

\textbf{Single architecture family.} Results demonstrated on GPT-2 variants only. Different architectures may exhibit different sensitivities.

\textbf{Single dataset family.} Experiments used RedPajama-v2. Domain-specific corpora may have different optimal thresholds.

\textbf{Deduplication dimension not fully validated.} The H-M2 deduplication analysis was inconclusive due to infrastructure issues.

\textbf{Single-parameter sweeps.} Perplexity and deduplication studied independently. Interaction effects remain unexplored.

\textbf{Synthetic/mock evaluation.} Due to lm-eval-harness integration issues, benchmark results are from simulation or mock evaluation, not actual model predictions.

\subsubsection{6.3 Broader Impact}

\textbf{Efficiency gains without new techniques.} Findings suggest practitioners can extract more performance from existing pipelines by calibrating curation parameters.

\textbf{Democratization of optimization.} Scale transfer findings mean smaller organizations can conduct systematic optimization at affordable scales.

\textbf{Methodology contribution.} The fixed-token experimental design and dose-response analysis framework can be applied to other curation parameters.
